% ================================================================
\section{The Nature of Mathematical Objects: Ontology}
\label{sec:ontology}
% ================================================================

\subsection{Benacerraf's Dilemma and the Cathedral's Response}

Benacerraf~\cite{benacerraf1973} identified a fundamental tension in
the philosophy of mathematics: a satisfactory account must explain
both the \emph{semantics} of mathematical statements (what they are
about) and the \emph{epistemology} of mathematical knowledge (how we
know them). Platonism gives a natural semantics (mathematical
statements are about abstract objects) but a mysterious epistemology
(how do we access these objects?). Formalism gives a natural
epistemology (mathematical knowledge is knowledge of formal
derivability) but an impoverished semantics (mathematical statements
are not ``about'' anything).

The Cathedral offers a distinctive response to this dilemma. Its
semantic content is exceptionally rich: the physics
dictionary~\cite{cathedral-physics} identifies 120+ structural
correspondences between the proof's mathematical objects and physical
phenomena. These are not post hoc analogies imposed by the author;
they emerge from the mathematical structure itself. The Gram matrix
\emph{is} a two-point correlator because it satisfies the axioms of a
two-point correlator (real, symmetric, positive definite). The Mertens
function \emph{is} a magnetization because it is the partial sum of
a $\{-1, 0, +1\}$-valued function with the cancellation properties
of a spin system.

At the same time, the Cathedral's epistemological standing is
unusually clear: the correspondences are not claims requiring faith
but \emph{theorems} verified by a proof assistant. When the physics
paper says ``the Parseval Bridge is unitarity,'' this is not a
metaphor but a machine-checked identity between two mathematical
expressions.

\subsection{Four Positions on Mathematical Objects}

What does the Cathedral's structure tell us about the ontological
status of mathematical objects? The classical positions offer
different readings.

\subsubsection{Platonism}

The Platonist holds that mathematical objects exist independently
of human minds and that mathematical truths are discovered, not
invented. The Cathedral's physics dictionary is natural evidence
for Platonism: if the integers possess genuine physical structure
(mass gaps, phase transitions, spectral statistics), then they
seem to exist in a way that transcends any particular formal
representation.

The graduated gates support this reading. The axiom
\lean{remainder\_bound\_abstract} was \emph{discovered} to be
false---not by fiat, but by investigation. The replacement theorems
were \emph{discovered} to be true. The language of discovery
presupposes the existence of something to be discovered.

\subsubsection{Formalism}

The formalist holds that mathematics is a game played with symbols
according to rules, and that the symbols do not ``refer'' to
anything. The Cathedral, as a Lean~4 term, is the paradigmatic
formalist object: a syntactic structure that type-checks.

But the physics dictionary creates a problem for formalism. If
mathematical objects are meaningless symbols, why do they have
physical content? The formalist might respond that the ``physics''
is merely a heuristic that guides the proof search, with no
ontological significance. This response is difficult to maintain
in the face of 120+ independently verified structural
correspondences, many of which were \emph{predicted} by the
physics dictionary and \emph{then} proved in Lean.

\subsubsection{Structuralism}

Shapiro's ante rem structuralism~\cite{shapiro1997} holds that
mathematical objects are positions in structures, and that
mathematical truths are truths about structural relationships.
On this view, the natural numbers are not Zermelo ordinals or
von Neumann ordinals (Benacerraf's~\cite{benacerraf1965}
observation that both are equally valid representations), but
the abstract structure that all representations share.

The Cathedral is deeply congenial to structuralism. Its central
achievement is to reveal the \emph{structure} of the Riemann
Hypothesis---the web of implications, equivalences, and
dependencies that connect RH to concrete analytic inequalities.
The physics dictionary is, in structuralist terms, a demonstration
that the structure of arithmetic is \emph{the same structure}
as that of certain physical systems.

\subsubsection{Intuitionism}

The intuitionist holds that mathematical objects are mental
constructions and that a mathematical statement is true only if a
constructive proof exists. The Cathedral is not constructive---it
uses Classical.choice throughout, and the converse direction relies
on the completeness of $L^2(0,1)$, a non-constructive result.

Nevertheless, the intuitionist can find value in the graduated
gates: each graduation is a construction that replaces an
assumption with a derivation. The progression $56 \to 1$ is a
sequence of increasingly constructive approximations to the full
proof. The Crown Axiom itself is a decidable predicate for each
fixed $N$ (one can compute $v^TGv$ and compare it to
$1 + K/\ln N$), even though the universal quantification over
all $N$ is not decidable.

\subsection{The Physics Dictionary as Ontological Evidence}

The most philosophically provocative feature of the Cathedral is
not its formal verification but its physics dictionary. The
companion paper~\cite{cathedral-physics} documents 120+ structural
correspondences between the mathematical objects of the proof and
physical phenomena:

\begin{center}
\small
\begin{tabular}{@{}ll@{}}
\toprule
\textbf{Mathematical Object} & \textbf{Physical Dual} \\
\midrule
Gram matrix $G(j,k)$ & Two-point quantum correlator \\
$d_N^2 = 1 - b^T G^{-1} b$ & Ground state (vacuum) energy \\
$M(x) = \sum_{n \leq x} \mu(n)$ & Magnetization \\
Parseval identity & Unitarity (probability conservation) \\
$\lambda_{\min}(G_N) > 0$ & Mass gap \\
$\lambda_{\min} \sim c/\ln N$ & Mass gap closing (phase transition) \\
Eigenvalue GOE statistics & Quantum chaos (thermalization) \\
Ground state scarring & Composite anchor localization \\
Perron contour shift & Crossing a phase boundary \\
Borel--Carath\'eodory & Maximum entropy principle \\
\bottomrule
\end{tabular}
\end{center}

\noindent These correspondences are not analogies. They are
\emph{structural identities}: the mathematical objects satisfy the
same axioms, obey the same inequalities, and exhibit the same
limiting behavior as their physical counterparts. The Gram matrix
does not merely ``resemble'' a quantum correlator---it IS a quantum
correlator, in the precise sense that it satisfies reflection
positivity, has a spectral decomposition with physically meaningful
eigenvalues, and its eigenvectors exhibit the localization patterns
(IPR, participation ratio) characteristic of quantum systems.

The philosophical question is: what do these correspondences
\emph{mean}?

\subsection{Three Interpretations}

\subsubsection{The Instrumentalist View}

The correspondences are useful heuristics for discovering and
organizing mathematical results, but have no ontological
significance. The physics vocabulary (``vacuum energy,''
``mass gap,'' ``thermalization'') is a convenient language for
navigating a complex proof, nothing more.

This view is difficult to sustain because the physics dictionary
has \emph{predictive power}. The identification of the off-diagonal
Gram structure as a ``Born approximation'' led to the prediction
that each entry should be non-positive (energy dissipation), which
was then proved as a theorem (\lean{eratio\_entry\_nonpos}). The
identification of the Möbius function as a ``spin variable'' led
to the SUSY decomposition, which was then formalized and verified.
Mere heuristics do not consistently generate provable theorems.

\subsubsection{The Structural Realist View}

The correspondences reflect a genuine structural isomorphism between
the integers and physical systems. The integers ``know physics''
not because they are physical objects, but because they share the
same abstract structure. Number theory and quantum mechanics are
two instantiations of a common mathematical architecture, and the
Cathedral reveals this architecture.

This is a moderate form of Tegmark's~\cite{tegmark2008} Mathematical
Universe Hypothesis: not the strong claim that the physical universe
IS a mathematical structure, but the weaker claim that certain
mathematical structures (in particular, the multiplicative structure
of the integers) and certain physical structures (quantum field
theories) are isomorphic at the level of their formal axioms.

Wigner's~\cite{wigner1960} ``unreasonable effectiveness of
mathematics in the natural sciences'' is usually discussed in the
direction physics $\to$ mathematics (we use math to describe
physics). The Cathedral exhibits the reverse direction:
mathematics $\to$ physics. The proof of a number-theoretic
equivalence spontaneously generates a complete physical
vocabulary, without any physical input.

\subsubsection{The Deep Identity View}

The strongest reading: mathematics and physics are not two domains
with a structural correspondence, but \emph{the same domain}
viewed from two perspectives. The integers are not ``like'' a
quantum field; they ARE the simplest quantum field. The Riemann
Hypothesis is not ``analogous to'' the assertion that the vacuum
is stable; it IS that assertion, in the only quantum field that
can be exactly described.

The Cathedral provides evidence for this view but cannot prove it.
What it can prove is that the structural correspondences are
exact---not approximate, not asymptotic, but machine-verified
identities between mathematical expressions. Whether this exactness
reflects a deep identity or merely a deep analogy is a question
that formal verification alone cannot answer. It may require a
conceptual framework that does not yet exist.

\subsection{The Born Approximation as Test Case}

A concrete example illuminates the ontological question. The
off-diagonal Gram entry decomposes as:
\[
  G(j,k) = \underbrace{\frac{\gcd(j,k)^2}{2jk}}_{\text{vacuum}}
  + \underbrace{E_{\text{ratio}}(j,k)}_{\text{Born approx.}}
  - \underbrace{\text{cotangent sums}}_{\text{scattering}}
  - \underbrace{\text{contact}}_{\text{self-energy}}
\]

Each term has a name because it \emph{behaves} like the named
physical quantity:
\begin{itemize}
  \item The ``vacuum'' term depends only on $\gcd(j,k)$---the
    shared divisor structure, independent of the individual modes.
  \item The ``Born approximation'' $E_{\text{ratio}}(j,k) =
    \frac{j-k}{2jk}\ln(k/j)$ is always $\leq 0$
    (\lean{eratio\_entry\_nonpos}): it dissipates energy, exactly
    as the Born approximation does in scattering theory.
  \item The ``scattering'' terms are cotangent sums---oscillatory
    phase factors that encode the fine structure of the
    divisibility relation.
  \item The ``contact'' term is local (diagonal correction).
\end{itemize}

The Born approximation non-positivity is a structural theorem: it
follows from the elementary inequality $(a-b)\ln(b/a) \leq 0$,
which holds for all positive reals. No physics was assumed; the
physics \emph{emerged} from the mathematics. The energy dissipation
is a theorem about logarithms, not a postulate about nature.

The instrumentalist says: we call it ``Born approximation''
because the label is helpful; the label has no ontological content.
The structural realist says: we call it ``Born approximation''
because it is isomorphic to the Born approximation; the isomorphism
is the content. The deep identity theorist says: it IS the Born
approximation, in the unique quantum field generated by the
integers; the usual Born approximation in physical scattering
theory is a different instantiation of the same abstract structure.

We incline toward the structural realist position, while
acknowledging that the evidence---120+ verified correspondences,
zero known counterexamples---does not conclusively establish any
one interpretation. What it does establish is that the question
is not idle: the Cathedral has made it empirically tractable.
